% ============================================================
% CHAPITRE 6 — Sprint 3 : Offres et candidatures
% ============================================================
\chapter{Sprint 3 --- Offres de Recrutement et Candidatures}

\section{Introduction}

Ce sprint constitue le cœur métier de la plateforme \textbf{DAAM}.
Il permet au RH de créer, publier et clôturer des \textbf{offres}
rattachées à une agence (donc à une zone), puis d'ouvrir ces offres aux
\textbf{candidats}~: dépôt de CV, passage du QCM, scoring, et suivi RH
des candidatures.

\section{Objectifs du Sprint}



\begin{itemize}
  \item permettre au RH de créer une offre (poste, compétences, agence,
        responsable de recrutement, QCM)~;
  \item générer automatiquement une \textbf{référence interne}
        (\texttt{REF-AAAA-NNN})~;
  \item gérer le cycle de vie \texttt{DRAFT} / \texttt{PUBLISHED} /
        \texttt{CLOSED}~;
  \item publier l'offre sur l'espace candidat (\texttt{/jobs}) tout en
        restant limité aux zones affectées au RH~;
  \item permettre au candidat de postuler avec un CV et un QCM~;
  \item calculer un score QCM et un score de matching CV~;
  \item permettre au RH de lister, filtrer et traiter les candidatures
        (non retenu, désisté, reset QCM, fiche de suivi)~;
  \item interdire le double dépôt sur la même offre.
\end{itemize}


\section{Spécification et Analyse des Besoins}

\subsection{Sprint Backlog}

\begin{table}[htbp]
  \centering
  \scriptsize
  \caption{Sprint Backlog du Sprint 3}
  \label{tab:sprint3-backlog}
  \begin{tabular}{|p{2.0cm}|c|p{4.1cm}|c|p{3.2cm}|c|}
    \hline
    \textbf{Fonctionnalité} & \textbf{ID} & \textbf{Scénario utilisateur} &
    \textbf{ID Tâche} & \textbf{Tâches} & \textbf{Est.} \\
    \hline
    Créer une offre
    & US-08
    & En tant que RH, je crée une offre liée à une agence et à un QCM.
    & T1
    & API \texttt{POST /recruitments}
    & 2 j \\
    \cline{4-6}
    & & & T2 & Formulaire \texttt{/rh/recruitments/new} & 2 j \\
    \hline
    Publier / clôturer
    & US-09
    & En tant que RH, je passe l'offre en publié ou clôturé.
    & T3
    & Statut DRAFT / PUBLISHED / CLOSED
    & 1 j \\
    \hline
    Consulter les offres
    & US-10
    & En tant que candidat, je consulte les offres publiées.
    & T4
    & API publique + pages \texttt{/jobs}
    & 1,5 j \\
    \hline
    Postuler
    & US-11
    & En tant que candidat, je dépose mon CV et passe le QCM.
    & T5
    & Upload CV + \texttt{POST /applications}
    & 2,5 j \\
    \hline
    Scoring
    & US-12
    & Le système calcule note QCM et matching CV.
    & T6
    & Scoring QCM + analyse CV
    & 2 j \\
    \hline
    Traitement RH
    & US-13
    & En tant que RH, je consulte et statue les candidatures de ma zone.
    & T7
    & Page \texttt{/rh/candidates} + PATCH statut
    & 2 j \\
    \hline
  \end{tabular}
\end{table}


\subsection{Diagramme de cas d'utilisation}

\begin{figure}[H]
  \centering
  \reportfig{img/uc-sprint3.png}
  \caption{Diagramme de cas d'utilisation des Sprint 3}
  \label{fig:uc-sprint3}
\end{figure}


\section{Description textuelle des cas d'utilisation}

\subsection{Cas d'utilisation <<~Créer une offre de recrutement~>>}

\begin{table}[htbp]
  \centering
  \small
  \caption{Description textuelle --- Créer une offre de recrutement}
  \label{tab:uc-create-offer}
  \begin{tabular}{|p{3.0cm}|p{10.2cm}|}
    \hline
    \textbf{Titre} & Créer une offre de recrutement \\
    \hline
    \textbf{Acteur} & RH \\
    \hline
    \textbf{Précondition} & RH authentifié, au moins une agence dans une zone affectée, QCM existant. \\
    \hline
    \textbf{Postcondition} & Offre créée (référence interne générée), statut brouillon ou publié. \\
    \hline
    \textbf{Scénario nominal}
    &
    \begin{enumerate}[leftmargin=*,itemsep=0.15em,topsep=0.2em]
      \item Le RH ouvre \texttt{/rh/recruitments/new}.
      \item Il saisit le poste, les compétences, choisit l'agence, le
            responsable de recrutement et le QCM.
      \item Le frontend envoie \texttt{POST /api/recruitment/recruitments}.
      \item Le service copie \texttt{zoneId} depuis l'agence, vérifie le
            périmètre RH et génère \texttt{REF-AAAA-NNN}.
      \item L'offre apparaît dans la liste RH.
    \end{enumerate}
    \\
    \hline
    \textbf{Alternatives}
    & Agence hors zone, ou RH sans affectation~: création refusée. \\
    \hline
  \end{tabular}
\end{table}


\subsection{Cas d'utilisation <<~Publier / clôturer une offre~>>}

\begin{table}[htbp]
  \centering
  \small
  \caption{Description textuelle --- Publier / clôturer une offre}
  \label{tab:uc-publish-offer}
  \begin{tabular}{|p{3.0cm}|p{10.2cm}|}
    \hline
    \textbf{Titre} & Publier ou clôturer une offre \\
    \hline
    \textbf{Acteur} & RH \\
    \hline
    \textbf{Précondition} & L'offre existe dans une zone du RH. \\
    \hline
    \textbf{Postcondition} & Statut \texttt{PUBLISHED} (visible candidats) ou \texttt{CLOSED}. \\
    \hline
    \textbf{Scénario nominal}
    &
    \begin{enumerate}[leftmargin=*,itemsep=0.15em,topsep=0.2em]
      \item Le RH édite l'offre (\texttt{/rh/recruitments/:id/edit}).
      \item Il choisit le statut (brouillon, publié, clôturé).
      \item \texttt{PUT /api/recruitment/recruitments/\{id\}} met à jour le statut.
    \end{enumerate}
    \\
    \hline
  \end{tabular}
\end{table}


\subsection{Cas d'utilisation <<~Postuler à une offre~>>}

\begin{table}[htbp]
  \centering
  \small
  \caption{Description textuelle --- Postuler à une offre}
  \label{tab:uc-apply}
  \begin{tabular}{|p{3.0cm}|p{10.2cm}|}
    \hline
    \textbf{Titre} & Postuler à une offre \\
    \hline
    \textbf{Acteur} & Candidat (\texttt{ROLE\_USER}) \\
    \hline
    \textbf{Précondition} & Offre \texttt{PUBLISHED} avec QCM~; candidat authentifié. \\
    \hline
    \textbf{Postcondition} & Candidature \texttt{SUBMITTED} avec scores QCM et matching CV. \\
    \hline
    \textbf{Scénario nominal}
    &
    \begin{enumerate}[leftmargin=*,itemsep=0.15em,topsep=0.2em]
      \item Le candidat ouvre \texttt{/jobs/:id/apply}.
      \item Il téléverse son CV (\texttt{POST /upload/cv}).
      \item Il répond au QCM puis valide.
      \item \texttt{POST /applications} enregistre les réponses, calcule
            le score et lance l'analyse CV.
      \item Le candidat retrouve la candidature dans
            \texttt{/jobs/applications}.
    \end{enumerate}
    \\
    \hline
    \textbf{Alternatives}
    & Déjà postulé~: refus. Anti-triche (changement d'onglet)~: score QCM forcé à 0. \\
    \hline
  \end{tabular}
\end{table}


\subsection{Cas d'utilisation <<~Traiter une candidature~>>}

\begin{table}[htbp]
  \centering
  \small
  \caption{Description textuelle --- Traiter une candidature}
  \label{tab:uc-process-app}
  \begin{tabular}{|p{3.0cm}|p{10.2cm}|}
    \hline
    \textbf{Titre} & Traiter une candidature \\
    \hline
    \textbf{Acteur} & RH \\
    \hline
    \textbf{Précondition} & Candidature dans une zone affectée au RH. \\
    \hline
    \textbf{Postcondition} & Statut mis à jour et/ou QCM réinitialisé. \\
    \hline
    \textbf{Scénario nominal}
    &
    \begin{enumerate}[leftmargin=*,itemsep=0.15em,topsep=0.2em]
      \item Le RH ouvre \texttt{/rh/candidates}.
      \item Il consulte matching CV, score QCM et fiche candidat.
      \item Il peut marquer \textbf{Non retenu} ou \textbf{Désisté},
            relancer l'analyse CV, ou autoriser un \textbf{Reset QCM}.
    \end{enumerate}
    \\
    \hline
    \textbf{Note} & La décision \textbf{Retenu} (entretien) est détaillée au sprint suivant. \\
    \hline
  \end{tabular}
\end{table}


\subsection{Autres cas d'utilisation}

\begin{itemize}
  \item \textbf{Lister / modifier / supprimer une offre}~: actions RH
        dans \texttt{/rh/recruitments}.
  \item \textbf{Consulter mes candidatures}~: espace candidat.
  \item \textbf{Autoriser un nouveau QCM}~: le RH réinitialise le
        questionnaire~; le candidat le repasse.
  \item \textbf{Consultation Admin}~: listes en lecture
        (\texttt{/admin/recruitments}, \texttt{/admin/applications}).
\end{itemize}


\section{Conception}

\subsection{Diagramme de classes}

Le domaine relie \textbf{Recruitment} (offre) à \textbf{Company} /
\textbf{Zone} et à un QCM, puis \textbf{JobApplication} au candidat et à
l'offre, avec scores QCM et matching CV.

\begin{figure}[H]
  \centering
  \reportfig{img/class-sprint3.png}
  \caption{Diagramme de classes des Sprint 3}
  \label{fig:class-sprint3}
\end{figure}


\subsection{Diagrammes de séquence}

\subsubsection{Diagramme de séquence <<~Créer une offre~>>}

\begin{figure}[H]
  \centering
  \reportfig{img/seq-create-offer.png}
  \caption{Diagramme de séquence --- Créer une offre}
  \label{fig:seq-create-offer}
\end{figure}

\subsubsection{Diagramme de séquence <<~Postuler à une offre~>>}

\begin{figure}[H]
  \centering
  \reportfig{img/seq-apply.png}
  \caption{Diagramme de séquence --- Postuler à une offre}
  \label{fig:seq-apply}
\end{figure}

\subsubsection{Diagramme de séquence <<~Traiter une candidature~>>}

\begin{figure}[H]
  \centering
  \reportfig{img/seq-process-app.png}
  \caption{Diagramme de séquence --- Traiter une candidature}
  \label{fig:seq-process-app}
\end{figure}


\section{Réalisation}

\subsection{Interface de gestion des offres (RH)}

La liste RH affiche les offres de la zone (titre, agence, statut, cohort).
La création / édition se fait via un formulaire unique.

\begin{figure}[H]
  \centering
  \reportfig{img/ui-rh-recruitments.png}
  \caption{Interface de gestion des offres (RH)}
  \label{fig:ui-rh-recruitments}
\end{figure}

\begin{figure}[H]
  \centering
  \reportfig{img/ui-rh-recruitment-form.png}
  \caption{Interface de création / modification d'une offre}
  \label{fig:ui-rh-recruitment-form}
\end{figure}


\subsection{Interface candidat --- offres et candidature}

Le candidat parcourt les offres publiées, consulte le détail, puis postule
(CV + QCM).

\begin{figure}[H]
  \centering
  \reportfig{img/ui-jobs.png}
  \caption{Interface de consultation des offres (candidat)}
  \label{fig:ui-jobs}
\end{figure}

\begin{figure}[H]
  \centering
  \reportfig{img/ui-apply.png}
  \caption{Interface de candidature (CV et QCM)}
  \label{fig:ui-apply}
\end{figure}


\subsection{Interface RH --- candidatures}

La page candidatures classe les postulants (matching CV, score QCM) et
propose les actions de traitement.

\begin{figure}[H]
  \centering
  \reportfig{img/ui-rh-candidates.png}
  \caption{Interface de gestion des candidatures (RH)}
  \label{fig:ui-rh-candidates}
\end{figure}


\section{Conclusion}

Le Sprint~3 a livré le flux complet \textbf{offre $\rightarrow$
candidature $\rightarrow$ préselection}~: le RH publie une offre zonée, le
candidat postule avec CV et QCM, le système score, et le RH statue
(non retenu, désisté, reset QCM).

Le chapitre suivant présente le Sprint~4, consacré aux \textbf{entretiens,
décisions, embauche et pilotage RH} (tableaux de bord, exports).

\clearpage
